% Core results summary paragraph
% Suitable for the Results section, immediately after Table~1.

On the feasible-geometry benchmark (eight scenarios, ten random seeds, $80$ episodes per method), the no-prediction guidance policy achieved a success rate of $75.0\%$.  The parametric prediction category (CV and CA, merged because they produced identical results under constant-velocity target motion) achieved $62.5\%$.  A paired $t$-test between no-prediction and parametric prediction yielded $p = 0.0013$ with Cohen's $d = -0.373$, indicating a small but statistically significant degradation.

Adding neural predictors revealed a striking architecture-dependent effect.  The frozen LSTM predictor recovered the no-prediction baseline exactly ($75.0\%$ success, mean return $-7.17 \pm 355.84$), succeeding on every \texttt{crossing\_right} episode where the parametric predictors failed.  In contrast, the frozen GRU predictor replicated the parametric failure mode ($62.5\%$ success, mean return $-110.78 \pm 400.42$, $p = 0.0013$ vs baseline).  All methods universally failed on the two head-on medium/far candidate scenarios ($0\%$), suggesting geometric infeasibility outside the training distribution.  The entire performance hierarchy is therefore determined by a single scenario: \texttt{crossing\_right} separates the $75\%$ methods (No-Prediction, LSTM) from the $62.5\%$ methods (Parametric, GRU).
